% ID: SEL-D0023-OB-E02
% ESTADO: validada. S024, cotejada con el examen.
\textbf{a)} Al completar el cuadrado, $f(x)=2(x-9)^2+38$.
El mínimo se alcanza para $x=9$, es decir, $\boxed{9000\text{ kg}}$,
con coste de $\boxed{38000\text{ euros}}$.

\textbf{b)} $f'(x)=4x-36$ y $f'(7)=-8<0$. Cerca de 7000 kg,
aumentar la producción reduce el coste.

\textbf{c)} Para el dominio físico $x\ge0$ se representa el arco
de parábola convexo con vértice $(9,38)$. Como
$f(x)=200$ equivale a $2x(x-18)=0$, se alcanza ese coste
con $\boxed{0\text{ o }18000\text{ kg}}$.
\begin{center}
\setlength{\unitlength}{4mm}
\begin{picture}(25,17)(-3,-1)
\qbezier(0,0)(10.5,0)(21,0)\put(21.2,0){$x$}
\qbezier(0,0)(0,7.25)(0,14.5)\put(0,15){$f$}
\qbezier(0,10)(10,-8)(20,14)
\put(0,10.4){\makebox(0,0)[r]{$(0,200)$}}
\put(9,1.5){\makebox(0,0)[t]{$(9,38)$}}
\put(18.1,9.4){\makebox(0,0)[l]{$(18,200)$}}
\end{picture}
\end{center}
